\section{字符串哈希}

模 $2^{61}-1$ 的多项式哈希：\texttt{Hash(s)} 预处理串 $s$ 后，\texttt{hash(l,r)} 返回 $s[l..r]$（下标 $0$ 起始、闭区间）的哈希值，两子串哈希相等即可判定二者相同。预处理 $\mathcal{O}(n)$，单次查询 $\mathcal{O}(1)$。

\texttt{Hash::P} 为静态幂表，多个实例共享并随串长懒扩充。基数 \texttt{K} 固定，题目开放 hack 时改成值为 $\frac{M}{2} \sim M-2$ 之间随机值即可。

\begin{lstlisting}
struct Hash {
    using u64 = unsigned long long;
    using u128 = unsigned __int128;
    constexpr static u64 M = (1ULL << 61) - 1;
    constexpr static u64 K = 1234567890123456789ULL;
    static inline std::vector<u64> P;
    constexpr static u64 norm(u64 x){
        x = (x & M) + (x >> 61);
        return x >= M ? x - M : x;
    }
    constexpr static u64 mul(u64 x,u64 y){
        u128 z = static_cast<u128>(x) * y;
        return norm((z & M) + (z >> 61));
    }
    std::vector<u64> H;
    Hash(const std::string& s):H(s.size() + 1){
        while (P.size() <= s.size()){
            P.push_back(P.empty() ? 1 : mul(P.back(),K));
        }
        for (int i = 0;i < s.size();i++){
            H[i + 1] = norm(mul(K,H[i]) + s[i]);
        }
    }
    u64 operator()(int l,int r) const {
        return norm(M + H[r + 1] - mul(H[l],P[r - l + 1]));
    }
};
\end{lstlisting}
